\chapter{Dimension}

In this chapter, all rings are assumed to be commutative, and algebras are associative, commutative, and unital.

Let A be a ring. If a is an ideal of A and n is a negative integer, we set \(a^{n} = A\) . For every prime ideal p of A, we denote by \(\kappa(\mathfrak{p})\) the residue field of the local ring \(A_{p}\) ; it is canonically isomorphic to the field of fractions of the ring A/p (II, § 3, no. 1, prop. 2). If A is local, we denote by \(m_{A}\) its maximal ideal, and \(\kappa_{A} = A/m_{A} = \kappa(m_{A})\) its residue field.

Let \(\rho: A \to B\) be a ring homomorphism. We denote by \(^a\rho\) the continuous map from \(\operatorname{Spec}(B)\) to \(\operatorname{Spec}(A)\) associated with \(\rho\) (II, § 4, no. 3). We say (V, § 2, no. 1, def. 1) that a prime ideal q of B lies over the prime ideal p of A if p is the image of q under \(^a\rho\) , that is, if \(\rho^{-1}(q) = p\) . Let M be a B-module; when we consider M as an A-module, it is always the A-module structure \(\rho_*(M)\) defined by the external law \((a, m) \mapsto \rho(a).m\) (A, II, p. 30).

Let M be an A-module. If U (resp. V) is an additive subgroup of A (resp. M), we denote by UV or U.V the additive subgroup of M generated by the products uv, for \(u \in U\) and \(v \in V\) . If S is a subset of A, we denote by SM the submodule \(\sum_{s \in S} sM\) of M; the ideal of A generated by S is equal to SA and we have S.M = SM.

